\subsection{A coadjoint orbit example}\label{theexample}
Consider the space $\Hom(\C^m,\C^n)$ of complex $n\times m$ matrices $Q$. The dual space is identified with matrices $P$ in $\Hom(\C^n,\C^m)$ via the trace pairing $\Tr(PQ)$. The holomorphic cotangent bundle $T^*_{1,0}\Hom(\C^m,\C^n)$ is then identified with the set of such pairs $(Q,P)$. The holomorphic symplectic form is
\begin{equation*}%\label{symp_form_on_v_space}
	\Omega((Q_1,P_1),(Q_2,P_2))=\Tr(P_2Q_1-P_1Q_2).
\end{equation*}
The groups ${\rm GL}_m\C$ and ${\rm GL}_n\C$ act symplectically by $\bigl(Qg^{-1},gP\bigr)$ and $\bigl(hQ,Ph^{-1}\bigr)$ with momentum maps
\begin{equation*}	%\label{momentum_maps}
	\mu_m(Q,P)=PQ,\qquad\text{and}\qquad\mu_n(Q,P)=QP.
\end{equation*}
Here we have identified $\mathfrak{gl}_n$ and $\mathfrak{gl}_m$ with their respective duals using the trace form.

Now suppose $m\le n$ and let $M$ be the open subset of the cotangent bundle where both $Q$ and $P$ have maximal rank. One can show (see \cite{dualpairs}) that $\mu_n$ is an orbit-map for the ${\rm GL}_m\C$-action on $M$ and that the fibres of $\mu_m$ are orbits of ${\rm GL}_n\C$. As the momentum map is Poisson we can identify the Poisson reduced space $M/{\rm GL}_m\C$ with $\mu_n(M)\subset\mathfrak{gl}_n\C^*$. For $\zeta\in\C^\times$, the orbit-reduced space $\mu_m^{-1}(\zeta\cdot\Id_m)/{\rm GL}_m\C$ is identified with the coadjoint orbit
\begin{equation}
	\label{coad_orbit}
	{\rm Orb}(\zeta)=\{\xi=QP\mid (Q,P)\in M,\,PQ=\zeta\cdot\Id_m\}\subset\mathfrak{gl}_n\C^*.
\end{equation}

Geometrically, $\xi$ determines a decomposition $\C^n=\Sigma^{(m)}\oplus \Lambda^{(n-m)}$, where ${\Sigma=\Imag Q}$ and $\Lambda=\ker P$. Conversely, such a decomposition determines an element $\xi$ where $\ker\xi=\Lambda$ and ${\xi|_\Sigma=\zeta\cdot\Id_\Sigma}$. Naturally, we see that the orbit is diffeomorphic to the symmetric space $\D_\C$ given~by
\begin{Definition}
	For $\mathbb{F}=\R,\C,\Q$, the \emph{space of decompositions} $\D_\mathbb{F}$ is the set
	\begin{equation}
		\label{elegant_A+B}
		\D_\mathbb{F}=(\Gr_\mathbb{F}(m;n)\times \Gr_\mathbb{F}(n-m;n))\setminus\Delta,
	\end{equation}
	where $\Delta$ is the closed subset of pairs $(\Sigma,\Lambda)$ with $\Sigma\cap\Lambda\ne\{0\}$.
\end{Definition}
\begin{Remark}\label{elegant_rmk}
	The symplectic form on $\D_{\C}$ admits a rather elegant description. Recall that tangent vectors to an element $\Sigma$ of the Grassmannian may be identified with linear maps ${\Sigma\rightarrow\C^n/\Sigma}$. Since $(\Sigma,\Lambda)$ determines a decomposition of $\C^n$, the tangent space to $\D_{\C}$ at this point may be identified with pairs of maps $(\alpha\colon\Sigma\rightarrow\Lambda,\,\beta\colon \Lambda\rightarrow\Sigma)$. The KKS form on ${\rm Orb}(\zeta)$ can be shown~to~be
	\[
		\Omega_{\rm KKS}((\alpha_1,\beta_1),(\alpha_2,\beta_2))=\zeta\Tr(\beta_2\alpha_1-\beta_1\alpha_2).
	\]
\end{Remark}
